% Generated draft: requires mathematical and source review.
\documentclass[11pt]{article}
\usepackage[margin=25mm]{geometry}
\usepackage{fontspec}
\setmainfont{Latin Modern Roman}
\usepackage{amsmath,amssymb,mathtools,unicode-math}
\setmathfont{Latin Modern Math}
\usepackage{enumitem,needspace,xurl,hyperref}
\setlength{\parindent}{0pt}
\setlength{\parskip}{5pt}
\setlength{\emergencystretch}{3em}
\newcommand{\N}{\mathbb{N}}
\newcommand{\Z}{\mathbb{Z}}
\newcommand{\Q}{\mathbb{Q}}
\newcommand{\R}{\mathbb{R}}
\newcommand{\C}{\mathbb{C}}
\newcommand{\F}{\mathbb{F}}
\newcommand{\A}{\mathbb{A}}
\newcommand{\PP}{\mathbb{P}}
\newcommand{\E}{\mathbb{E}}
\newcommand{\eps}{\varepsilon}
\newcommand{\dd}{\,\mathrm{d}}
\begin{document}
\section*{UnsolvedMath ID: 30003124}
\textbf{Source problem code:} catalogue entry
\subsection{Definitions and mathematical statement}
Under the standard interpretation of the supplied notation, let \(p\) be a fixed prime and let \(q\) tend to infinity through the positive integers \(\Z_{>0}\). These conventions are implicit in the source statement. For \(q\in\Z_{>0}\), define
\[
v_p(q)=\max\{k\in\Z_{\geq 0}:p^k\mid q\},\qquad |q|_p=p^{-v_p(q)}.
\]
For \(t\in\R\), define its distance to the nearest integer by
\[
\|t\|=\min_{m\in\Z}|t-m|.
\]
The \(p\)-adic Littlewood conjecture referred to in the record asserts that, for every prime \(p\) and every \(x\in\R\),
\[
\liminf_{\substack{q\to\infty\\q\in\Z_{>0}}}q|q|_p\|qx\|=0.
\]
For a fixed prime \(p\), define the set of hypothetical counterexamples by
\[
\mathcal C_p=\left\{x\in\R:
\liminf_{\substack{q\to\infty\\q\in\Z_{>0}}}q|q|_p\|qx\|>0\right\}.
\]
Equivalently, \(x\in\mathcal C_p\) means that there exist \(c>0\) and \(Q\in\Z_{>0}\) such that
\[
q|q|_p\|qx\|\geq c\qquad\text{for every integer }q\geq Q.
\]
\textbf{Problem.} Determine necessary structural properties of a real number \(x\) under the hypothesis \(x\in\mathcal C_p\). The source asks how such a number should behave, but specifies neither a particular property to prove nor a precise classification criterion. Thus this is a conditional structural research question; it does not assert that \(\mathcal C_p\) is nonempty. The intended scope of ``behave'' remains unresolved.
\subsection{Short English statement}
For a fixed prime \(p\), investigate the structure of a real number that would violate the \(p\)-adic Littlewood conjecture: its quantities \(q|q|_p\|qx\|\) must stay bounded away from zero for all sufficiently large positive integers \(q\). The source does not specify which further structural properties are sought or whether a complete characterization is intended.
\subsection{Sources}
\begin{itemize}
\item Supplied local UnsolvedMath registry record, permanent ID \(30003124\), title: \emph{Structure of Counterexamples to the p-Adic Littlewood Conjecture}; exact source problem code: \texttt{catalogue entry}. Collection: Oberwolfach Reports Open Problems. No proposer is identified; the supplied proposed year is 2016.
\item Supplied report citation: \emph{Diophantische Approximationen (2017)}, Oberwolfach Reports, DOI 10.4171/owr/2016/21, pp.\ 1099--1187. \url{https://doi.org/10.4171/owr/2016/21}. The background describes this collection as extracted from workshop reports and their open-problem discussions. The supplied citation's 2017 label and its 2016 DOI component are retained without resolving their relationship.
\item Badziahin, \emph{Continued Fractions of Potential Counterexamples}. \url{https://arxiv.org/abs/1406.3594}.
\item Badziahin, \emph{Estimating the Lower Limit in the p-Adic Littlewood Conjecture}. \url{https://doi.org/10.1080/10586458.2025.2595008}.
\item Supplied literature assessment, checked on \texttt{2026-08-21}: the record labels the problem open and states that the conjecture remains unresolved, with the literature deriving strong structural restrictions on hypothetical counterexamples. This is supplied dated curation derived from the OWR material, not an independently verified literature assessment; time-sensitive status and cited claims require independent verification.
\item The registry labels the statement \texttt{corrected\_verified} and reports that its mathematical content was preserved while notation, typography, or omitted definitions were repaired for self-contained publication. These are supplied editorial assertions, not independent mathematical verification of this draft.
\end{itemize}
\end{document}
